% artikkel_VII_fitnesslandskap.tex
% Artikkel VII: Det empiriske fitnesslandskapet
% STOLAR-prosjektet - Kvantitativ designhistorie
% Spraak: Nynorsk

\documentclass[10pt, a4paper, twocolumn]{article}

\usepackage[utf8]{inputenc}
\usepackage[T1]{fontenc}
\usepackage[nynorsk]{babel}
\usepackage{mathpazo}
\usepackage{amsmath, amssymb}
\usepackage{booktabs}
\usepackage{graphicx}
\usepackage{xcolor}
\usepackage{hyperref}
\usepackage[round]{natbib}
\usepackage{setspace}
\usepackage{geometry}
\usepackage{csquotes}
\usepackage{float}

\geometry{left=1.0cm, right=1.0cm, top=1.8cm, bottom=1.8cm, columnsep=0.6cm}
\setstretch{1.15}
\graphicspath{{fig/}}

\title{%
  \textbf{Det empiriske fitnesslandskapet}\\[0.4em]
  \large Fr\r{a}a metafor til topografi i 2\,284 stolar%
}
\author{%
  Iver Raknes Finne\\
  \small Arkitektur- og designh\o gskolen i Oslo (AHO)\\
  \small \texttt{iver.raknes.finne@stud.aho.no}
}
\date{}

\begin{document}
\maketitle

% ================================================================
\begin{abstract}
\noindent Fitnesslandskapet har vore designteoriens mest lovande metafor sidan \citet{wright1932} og \citet{kauffman1993} gav det matematisk form. Likevel har ingen studie konstruert landskapet som eit empirisk objekt for formgjeving. Denne artikkelen gjer nettopp det. Med utgangspunkt i STOLAR-databasen ($N = 2\,284$ stolar, 1280--2024) operasjonaliserer me det fleirdimensjonale fitnesslandskapet gjennom fem kvantifiserbare dimensjonar: materialentropi ($H'$), teknikkompleksitet ($T_c$), dimensjonell varians ($\sigma_D$), geografisk spreiing ($G_s$) og estimert materiell intensitet ($I_m$). Me konstruerer toppunktprofilar for 13 europeiske stilperiodar og viser at kvar stil er eit lokalt optimum med distinkt topografi. Postmodernismen produserer det mest ekstreme toppunktet (intensitetsindeks $\mathcal{I} = 182{,}6$), driven av det h\o gaste materielle investeringsniv\r{a}et i datasettet (gjennomsnittleg vekt 41{,}9 kg); barokken er det nest h\o gaste ($\mathcal{I} = 129{,}2$). Empire produserer det strammaste toppunktet ($\sigma_D = 0{,}066$), og modernismen det breiaste plataaet (breiddeindeks $\mathcal{B} = 32{,}1$). Landskapsskifte, punkta der heile topografien endrar seg, vert identifiserte gjennom Jaccard-avstandar mellom paaf\o lgjande stilars material- og teknikkrepertoar. Dei tre st\o rste skifta er jugend til funksjonalisme ($\hat{J} = 0{,}853$), historisme til jugend ($\hat{J} = 0{,}741$), og empire til historisme ($\hat{J} = 0{,}724$). Desse skifta samanfell med kjende materialrevolusjonar: industrialiseringa, handverksreaksjonen og st\r{a}lr\o yr-innovasjonen. Resultata viser at designhistoria ikkje er ein monoton marsj mot betre funksjon, men ein \emph{vandring over eit ruglete terreng} der lokale optimum vert destabiliserte av materielle og teknologiske innovasjonar.

\medskip
\noindent\textbf{N\o kkelord:} fitnesslandskap, topografi, stilperiode, Jaccard-avstand, materialentropi, seleksjonstrykk, kvantitativ designhistorie, Form Follows Fitness
\end{abstract}

\vspace{1em}
\hrule
\vspace{1.5em}


% ================================================================
\section{Innleiing: behovet for eit empirisk landskap}

Designteori har lenge brukt biologiske metaforar. \citet{sullivan1896} henta ``form follows function'' fr\r{a}a organisk naturfilosofi. \citet{steadman2008} skreiv heile b\o ker om evolution\ae r analogi i arkitektur. \citet{kauffman1993} formaliserte fitnesslandskapet matematisk, og \citet{frenken2006} demonstrerte det for teknologisk innovasjon. Men ingen av desse overf\o rte metaforen til \emph{empirisk m\r{a}albar topografi} for formgjevne objekt.

Denne mangelen er eit problem. Utan empirisk substans er fitnesslandskapet berre ei pedagogisk illustrasjon, ei skisse p\r{a}a ein tavle. \citet{kauffman1993} sine $NK$-modellar har presise matematiske eigenskapar: talet p\r{a}a lokale optima veks eksponentielt med koplingsgraden $K$, og ``kompleksitetskatastrofen'' gjer global optimalisering umogleg i sterkt kopla system. \citet{gavrilets2004} utvida modellen med ``holey landscapes'' der n\o ytrale r\aa s (ridges) knyter saman lokale toppunkt. Men korleis ser dette ut for \emph{stolar}?

Artikkel~IV i denne serien definerte fitnesslandskapet formelt og demonstrerte gjennom ni kumulative bevis at form ikkje er dedusert fr\r{a}a funksjon, men selektert av eit fleiraksig tilpassingslandskap \citep{finne2026}. Tilpassingsfunksjonen $\Phi(\mathbf{x}) = \sum_i w_i \cdot f_i(\mathbf{x})$ vart definert med seks komponentar (funksjonell yting, materialtilgang, teknologisk kapasitet, produksjonskostnad, kulturell meining, geografisk kontekst). Men $\Phi$ var kvalitativt definert, ikkje kvantitativt estimert. Landskapet var postulert, ikkje kartlagt.

\begin{quote}
\textbf{G\r{a}atet i denne artikkelen:} Me operasjonaliserer fitnesslandskapet gjennom fem m\r{a}albare dimensjonar, konstruerer toppunktprofilar for 13 stilperiodar, og identifiserer landskapsskifte gjennom Jaccard-analyse. Resultatet er ikkje ein metafor, men eit kart.
\end{quote}

Artikkelen er strukturert slik: seksjon~2 definerer dei fem dimensjonane og koplar kvar til ei spesifikk datakjelde i STOLAR-databasen. Seksjon~3 presenterer toppunktprofilane for dei 13 stilperiodane med mest enn 17 observasjonar. Seksjon~4 analyserer landskapsskifte gjennom Jaccard-avstandar mellom kronologisk tilgrensande stilar. Seksjon~5 dr\o ftar implikasjonar og avgrensingar. Seksjon~6 konkluderer.


% ================================================================
\section{Dei fem dimensjonane}

Fitnesslandskapet krev operasjonalisering: kva aksar definerer terrenget? Me vel fem dimensjonar som til saman fangar det materielle, teknologiske, geometriske, geografiske og intensive aspektet ved kvar stilperiode. Kvar dimensjon er utleidd fr\r{a}a eksisterande kolonnar i STOLAR-databasen og kan rekalkulerast av kven som helst med tilgang til datasettet.

Tabell~\ref{tab:operasjonalisering} gjev eit oversyn.

\begin{table}[H]
\centering
\caption{Operasjonalisering av dei fem landskapsdimensjonane.}
\label{tab:operasjonalisering}
\begin{tabular}{lccl}
\toprule
\textbf{Dimensjon} & \textbf{Symbol} & \textbf{Eining} & \textbf{Datakjelde} \\
\midrule
Materialentropi       & $H'$       & bits           & Materialar \\
Teknikkompleksitet   & $T_c$      & bits           & Teknikk \\
Dimensjonell varians & $\sigma_D$ & dimensjonslaus & H, W, D \\
Geografisk spreiing  & $G_s$      & indeks         & Nasjonalitet \\
Materiell intensitet & $I_m$      & $z$-skore      & Vekt, volum, mat.tal \\
\bottomrule
\end{tabular}
\end{table}


\subsection{Materialentropi ($H'$)}

\textbf{Materialentropien} m\r{a}aler kor mange materialar ein stilperiode nyttar og kor jamt dei er fordelte. Formelt er dette Shannon-entropien \citep{shannon1948} over materialpaletten innanfor kvar stilperiode:

\begin{equation}
  H' = -\sum_{i=1}^{S} p_i \log_2 p_i
\end{equation}

\noindent der $p_i$ er andelen av materiale $i$ blant alle materialforekomstar i stilperioden, og $S$ er talet p\r{a}a distinkte materialar. Artikkel~I demonstrerte at $H'$ for heile databasen stig fr\r{a}a 1{,}0 bits (1200-talet) til 5{,}07 bits (1900-talet). Her bereknar me $H'$ \emph{per stilperiode} snarare enn per hundre\r{a}ar.

Funksjonalismen har den h\o gaste materialentropien ($H' = 4{,}48$ bits) av dei 13 stilperiodane: 47 stolar nyttar eit breitt spekter av materialar fr\r{a}a st\r{a}al og kryssfiner til plast og aluminium. Hepplewhite har den l\r{a}agaste ($H' = 2{,}99$ bits) blant dei store stilane: 28 stolar dominert av mahogni og satintre. Skilnaden er ikkje tilfeldig: ho reflekterer den materielle konteksten kvar stil opererer i.


\subsection{Teknikkompleksitet ($T_c$)}

\textbf{Teknikkompleksiteten} m\r{a}aler repertoaret av konstruksjonsteknikkar innanfor kvar stilperiode, ogs\r{a}a som Shannon-entropi:

\begin{equation}
  T_c = -\sum_{j=1}^{T} q_j \log_2 q_j
\end{equation}

\noindent der $q_j$ er andelen av teknikk $j$. Artikkel~III viste at teknikk-entropien stig parallelt med materialentropien, fr\r{a}a 1{,}0 til 3{,}67 bits over 700 \r{a}ar. Funksjonalismen toppar ogs\r{a}a denne dimensjonen ($T_c = 3{,}09$ bits): han nyttar formboyging, sveising, sproytest\o yping, tapping og polstring samstundes. Hepplewhite er l\r{a}agast ($T_c = 1{,}30$ bits): eit smalt teknikkrepertoar dominert av tapping og polstring.


\subsection{Dimensjonell varians ($\sigma_D$)}

\textbf{Dimensjonell varians} m\r{a}aler kor stramt eller spreidd formrommet er innanfor ein stilperiode. Me definerer det som gjennomsnittleg variasjonskoeffisient (CV) over dei tre hovuddimensjonane:

\begin{equation}
  \sigma_D = \frac{1}{3}\left(\text{CV}_H + \text{CV}_W + \text{CV}_D\right)
\end{equation}

\noindent der $\text{CV}_X = s_X / \bar{X}$ er standardavviket delt p\r{a}a gjennomsnittet for dimensjon $X$.

Ein l\r{a}ag $\sigma_D$ tyder at stilperioden kanaliserer forma inn i eit smalt band: eit skarpt, h\o gt toppunkt i landskapet. Ein h\o g $\sigma_D$ tyder at stilperioden er eit breitt, flatt plat\r{a}a der formvariasjonen er stor.

Empire har den l\r{a}agaste $\sigma_D = 0{,}066$: empirestolane er n\ae rast identiske i proporsjonar ($\text{CV}_H = 0{,}013$, den strammaste profilen i heile datasettet). I kontrast har postmodernismen $\sigma_D = 0{,}351$: standardavviket for h\o gde er 41\,\% av gjennomsnittet. Postmodernismen er ikkje eit toppunkt; det er eit plat\r{a}a.


\subsection{Geografisk spreiing ($G_s$)}

\textbf{Geografisk spreiing} m\r{a}aler kor mange nasjonalitetar som er representerte innanfor ein stilperiode, justert for utvalsstorleik gjennom Margalef sin rikdomsindeks:

\begin{equation}
  G_s = \frac{S_{\text{nat}} - 1}{\ln N}
\end{equation}

\noindent der $S_{\text{nat}}$ er talet p\r{a}a distinkte nasjonalitetar og $N$ er talet p\r{a}a stolar med registrert nasjonalitet.

Rokokko har den h\o gaste geografiske spreiinga ($G_s = 1{,}82$): 27 stolar fr\r{a}a mange europeiske land. Rokokko var ein genuin pan-europeisk stil. Jugend har den l\r{a}agaste ($G_s = 0{,}00$): alle 21 jugendstolane i datasettet kjem fr\r{a}a same nasjonalitetskategori, noko som reflekterer jugend sin sterkt nasjonale karakter (norsk dragestil, finsk jugend, osb. er registrerte under eitt land i dette delsettet). Modernismen har h\o g spreiing ($G_s = 1{,}44$): internasjonalismen er ikkje berre ein ideologi, men ein m\r{a}albar eigenskap.


\subsection{Estimert materiell intensitet ($I_m$)}

\textbf{Materiell intensitet} er ein samansett indeks som fangar kor mykje materiale ein stilperiode investerer i kvar stol. Me definerer han som summen av $z$-skoarar for tre variablar:

\begin{equation}
  I_m = z(\bar{W}) + z(\bar{M}_c) + z(\bar{V})
\end{equation}

\noindent der $\bar{W}$ er gjennomsnittleg estimert vekt (kg), $\bar{M}_c$ er gjennomsnittleg materialtal per stol, og $\bar{V}$ er gjennomsnittleg begrensingsvolum ($H \times W \times D$, cm$^3$). Kvar $z$-skoar er standardisert p\r{a}a tvers av dei 13 stilperiodane.

Postmodernismen dominerer denne dimensjonen ($I_m = 6{,}68$): gjennomsnittleg vekt er 41{,}9 kg og gjennomsnittleg begrensingsvolum er 617\,313 cm$^3$. Desse tala er drivne av skulpturelle stolar som Sottsass sin Carlton (1981) og liknande utstillingsobjekt. Empire ($I_m = 1{,}07$) og Hepplewhite ($I_m = 0{,}95$) har ogs\r{a}a positiv materiell intensitet, driven av tunge mahognistolar. Funksjonalismen er mest materialsparsom ($I_m = -1{,}70$): lettare stolar i st\r{a}al og kryssfiner.


% ================================================================
\section{Toppunktprofilane: 13 stilperiodar som lokale optimum}

Kvar stilperiode i STOLAR-databasen med minst 17 observasjonar f\r{a}ar ein \textbf{toppunktprofil}: ein femdimensjonal vektor som beskriv posisjonen i fitnesslandskapet:

\begin{equation}
  \mathbf{p}_s = (H'_s,\; T_{c,s},\; \sigma_{D,s},\; G_{s,s},\; I_{m,s})
\end{equation}

Tabell~\ref{tab:toppprofilar} presenterer dei 13 profilane i kronologisk rekkjefolge. To samansette indeksar oppsummerer kvar profil:

\textbf{Intensitetsindeksen} $\mathcal{I}$ m\r{a}aler den samla ``h\o gda'' av toppunktet, definert som den euklidske norma av den min-maks-normaliserte profilen, skalert med 100:

\begin{equation}
  \mathcal{I}_s = 100 \cdot \left\|\hat{\mathbf{p}}_s\right\| = 100 \cdot \sqrt{\sum_{d=1}^{5} \hat{p}_{s,d}^2}
\end{equation}

\textbf{Breiddeindeksen} $\mathcal{B}$ m\r{a}aler kor ``flatt'' eller ``breitt'' toppunktet er, definert som produktet av dimensjonell varians og geografisk spreiing, skalert med 100:

\begin{equation}
  \mathcal{B}_s = 100 \cdot \sigma_{D,s} \cdot G_{s,s}
\end{equation}

Ein stil med h\o g $\mathcal{I}$ og l\r{a}ag $\mathcal{B}$ er eit \emph{skarpt, h\o gt toppunkt}: konsentrert materiell investering i eit smalt formrom. Ein stil med h\o g $\mathcal{B}$ er eit \emph{breitt plat\r{a}a}: spreidd over eit vidt formrom og mange land.


\begin{table*}[!ht]
\centering
\caption{Toppunktprofilar for 13 stilperiodar ($N \geq 17$). $H'$: materialentropi (bits). $T_c$: teknikkompleksitet (bits). $\sigma_D$: dimensjonell varians. $G_s$: geografisk spreiing (Margalef). $I_m$: materiell intensitet ($z$-skoar). $\mathcal{I}$: intensitetsindeks. $\mathcal{B}$: breiddeindeks.}
\label{tab:toppprofilar}
\begin{tabular}{lrrrrrrrrr}
\toprule
\textbf{Stil} & $N$ & $H'$ & $T_c$ & $\sigma_D$ & $G_s$ & $I_m$ & $\mathcal{I}$ & $\mathcal{B}$ \\
\midrule
Renessanse       & 18 & 3,36 & 2,22 & 0,269 & 0,910 & 0,25  & 116,0 & 24,52 \\
Barokk           & 92 & 4,21 & 2,27 & 0,181 & 1,077 & 0,59  & 129,2 & 19,49 \\
R\'{e}gence      & 27 & 3,12 & 2,27 & 0,096 & 0,390 & 0,14  &  75,9 &  3,76 \\
Rokokko          & 27 & 2,63 & 1,83 & 0,081 & 1,820 & $-$1,60 & 106,4 & 14,78 \\
Hepplewhite      & 28 & 2,99 & 1,30 & 0,064 & 0,481 & 0,95  &  54,4 &  3,06 \\
Louis XVI        & 22 & 2,71 & 1,64 & 0,146 & 0,558 & 0,02  &  55,7 &  8,14 \\
Nyklassisisme    & 26 & 3,69 & 2,43 & 0,079 & 0,417 & $-$1,16 &  94,3 &  3,27 \\
Empire           & 36 & 3,14 & 2,27 & 0,066 & 0,805 & 1,07  &  88,1 &  5,28 \\
Historisme       & 64 & 3,92 & 2,72 & 0,133 & 0,638 & $-$1,48 & 117,9 &  8,49 \\
Jugend           & 21 & 2,94 & 1,79 & 0,122 & 0,000 & $-$0,54 &  49,9 &  0,00 \\
Funksjonalisme   & 47 & 4,48 & 3,09 & 0,220 & 1,038 & $-$1,70 & 161,8 & 22,79 \\
Modernisme       & 17 & 3,70 & 3,00 & 0,222 & 1,443 & $-$1,64 & 150,8 & 32,06 \\
Postmodernisme   & 27 & 3,70 & 2,83 & 0,351 & 0,739 & 6,68  & 182,6 & 25,94 \\
\bottomrule
\end{tabular}
\end{table*}


\subsection{Det mest ekstreme toppunktet: postmodernisme ($\mathcal{I} = 182{,}6$)}

Postmodernismen produserer den h\o gaste intensitetsindeksen i heile datasettet. Drivkrafta er den materielle intensiteten: gjennomsnittleg vekt er 41{,}9 kg, nesten dobbelt s\r{a}a mykje som barokken (23{,}3 kg) og seks gonger s\r{a}a mykje som funksjonalismen (12{,}3 kg). Gjennomsnittleg begrensingsvolum er 617\,313 cm$^3$, meir enn dobbelt av barokken (295\,084 cm$^3$).

Desse tala reflekterer postmodernismen sin karakter som eit estetisk oppr\o r: stolar som Ettore Sottsass sin Carlton (1981), Philippe Starck sine skulpturelle uttrykk og Ron Arad sine st\r{a}alskulpturar er ikkje optimerte for funksjon eller \o konomi, men for kulturell provokasjon. Det materielle investeringsniv\r{a}aet er h\o gast \emph{nettopp fordi} funksjon er minst avgjerande. Postmodernismen er empirisk bevis for FFF sin kjernep\r{a}astand: n\r{a}ar funksjonstrykket er svakt, eksploderer formrommet.

Samstundes er postmodernismen det mest \emph{spreidde} toppunktet ($\sigma_D = 0{,}351$): standardavviket for h\o gde er 41\,\% av gjennomsnittet. Det er eit toppunkt som er b\r{a}ade h\o gt og breitt: ein vulkan, ikkje ein n\r{a}al.


\subsection{Det nest h\o gaste: barokk ($\mathcal{I} = 129{,}2$)}

Barokken er den nest mest intense stilperioden. Materialentropien er h\o g ($H' = 4{,}21$ bits): eik, mahogni, silke, messing, forgylling og lakkering inng\r{a}ar alle i repertoaret. Den geografiske spreiinga er betydeleg ($G_s = 1{,}08$): barokken er ein pan-europeisk stil med produksjon i mange land. Gjennomsnittleg vekt er 23{,}3 kg, den h\o gaste blant dei pre-modernistiske stilane.

I motsetnad til postmodernismen er barokken meir \emph{fokusert}: $\sigma_D = 0{,}181$ tyder at formrommet er konsentrert rundt ein bestemt konfigurasjon (h\o g, smal, tung). Barokken er ein fjelltopp med bratte sider: eit distinkt lokalt optimum der materiell rikdom, dreia ornamentikk og vertikal autoritet konvergerer. Funksjonalismen ($\mathcal{I} = 161{,}8$) ligg mellom dei to, driven av h\o gast materialentropi ($H' = 4{,}48$) og teknikkompleksitet ($T_c = 3{,}09$), men med l\r{a}ag materiell intensitet ($I_m = -1{,}70$): mykje materiell diversitet, men lite materiell investering per stol. Funksjonalismen er eit \emph{breitt toppunkt} ($\mathcal{B} = 22{,}8$): mange materialar, mange teknikkar, stor formvariasjon.


\subsection{Det strammaste toppunktet: empire ($\sigma_D = 0{,}066$)}

Empire har den l\r{a}agaste dimensjonelle variansen av alle 13 stilar (saman med Hepplewhite, $\sigma_D = 0{,}064$). Empire-stolar er n\ae rast identiske: $\text{CV}_H = 0{,}013$, gjennomsnittleg h\o gde 80{,}8 cm med eit standardavvik p\r{a}a berre 1{,}1 cm. 36 stolar, alle i mahogni, alle med tapping og polstring, alle med same nyklassisistiske proporsjonar.

Denne ekstremt l\r{a}age variansen er ei kvantitativ manifestering av det Artikkel~I kalla \textbf{mahogniens boge}: det koloniale materialmonopolet kanaliserer forma s\r{a}a sterkt at dimensjonell variasjon kollapsar mot null. Empire er det skarpaste toppunktet i fitnesslandskapet: ein n\r{a}al, ikkje ein vulkan. Kanaliseringskreftene (materiale, teknikk, proporsjonslare) er alle maksimalt synkroniserte.

Breiddeindeksen stadfestar biletet: $\mathcal{B} = 5{,}28$ for empire mot $\mathcal{B} = 32{,}1$ for modernisme. Empire er seks gonger smalare enn modernisme i landskapstopografien.


\subsection{Mellomstilar: navigasjon mellom ytterpunkt}

Mellom dei mest ekstreme profilane ligg stilar som representerer mellomposisjonar i landskapet. R\'{e}gence ($\mathcal{I} = 75{,}9$, $\mathcal{B} = 3{,}8$) er eit moderat, konsentrert toppunkt med strame proporsjonar (114{,}5 cm gjennomsnittleg h\o gde, $\text{CV}_H = 0{,}066$) og avgrensa materialpalett. Historisme ($\mathcal{I} = 117{,}9$, $\mathcal{B} = 8{,}5$) er eit breitt, moderat h\o gt toppunkt: 64 stolar med h\o g materialentropi ($H' = 3{,}92$) men negativt materiell intensitet ($I_m = -1{,}48$), noko som reflekterer industriell masseproduksjon med rimelegare materialar.

Nyklassisisme ($\mathcal{I} = 94{,}3$) og rokokko ($\mathcal{I} = 106{,}4$) er n\ae rast spegelvende: nyklassisismen har h\o g materialentropi ($H' = 3{,}69$) men l\r{a}ag geografisk spreiing ($G_s = 0{,}42$), medan rokokko har l\r{a}ag materialentropi ($H' = 2{,}63$) men h\o gast geografisk spreiing ($G_s = 1{,}82$). Dei okkuperer ulike regionar av same h\o gdeband i landskapet: ein materialtett lokal stil mot ein materielt enkel internasjonal stil.


% ================================================================
\section{Landskapsskifte: Jaccard-avstandar mellom stilar}

Eit fitnesslandskap er ikkje statisk. N\r{a}ar ein ny materialteknologi vert tilgjengeleg, eller n\r{a}ar ein handelsveg opnar eller kollapsar, endrar heile topografien seg. Tidlegare toppunkt kan verte dalar; nye toppunkt kan oppst\r{a}a. Desse overgangane kallar me \textbf{landskapsskifte}.

Me kvantifiserer landskapsskifte gjennom \textbf{Jaccard-avstanden} \citep{magurran2004} mellom materialpalettane til kronologisk tilgrensande stilperiodar:

\begin{equation}
  \hat{J}(s_1, s_2) = 1 - \frac{|M_{s_1} \cap M_{s_2}|}{|M_{s_1} \cup M_{s_2}|}
\end{equation}

\noindent der $M_{s}$ er mengda av distinkte materialar brukt i stilperiode $s$. Ein h\o g $\hat{J}$ tyder at dei to stilane nyttar radikalt ulike materialar: landskapet har endra seg mellom dei. Ein l\r{a}ag $\hat{J}$ tyder at materialpaletten er stabil: stilane navigerer innanfor same landskap.

Tabell~\ref{tab:jaccard} presenterer Jaccard-avstandar for b\r{a}ade material- og teknikkrepertoar mellom alle kronologisk tilgrensande stilpar.


\begin{table}[H]
\centering
\caption{Jaccard-avstandar mellom kronologisk tilgrensande stilperiodar. $\hat{J}_{\text{mat}}$: materialpalett. $\hat{J}_{\text{tekn}}$: teknikkrepertoar. $\Delta H'$: endring i materialentropi.}
\label{tab:jaccard}
\begin{tabular}{lrrr}
\toprule
\textbf{Stilpar} & $\hat{J}_{\text{mat}}$ & $\hat{J}_{\text{tekn}}$ & $\Delta H'$ \\
\midrule
Ren. $\rightarrow$ Barokk         & 0,607 & 0,667 & 0,85 \\
Barokk $\rightarrow$ R\'{e}g.     & 0,481 & 0,300 & 1,09 \\
R\'{e}g. $\rightarrow$ Rokokko    & 0,500 & 0,500 & 0,49 \\
Rokok. $\rightarrow$ Heppl.       & 0,667 & 0,714 & 0,35 \\
Heppl. $\rightarrow$ Louis XVI    & 0,667 & 0,667 & 0,27 \\
L.\,XVI $\rightarrow$ Nyklass.    & 0,556 & 0,667 & 0,97 \\
Nyklass. $\rightarrow$ Empire     & 0,421 & 0,400 & 0,54 \\
Empire $\rightarrow$ Hist.        & 0,724 & 0,417 & 0,77 \\
Hist. $\rightarrow$ Jugend        & 0,741 & 0,667 & 0,98 \\
Jugend $\rightarrow$ Funkj.       & 0,853 & 0,750 & 1,53 \\
Funkj. $\rightarrow$ Modern.      & 0,722 & 0,333 & 0,78 \\
Mod. $\rightarrow$ Postmod.       & 0,667 & 0,417 & 0,00 \\
\bottomrule
\end{tabular}
\end{table}


\subsection{Dei tre st\o rste landskapsskifta}

Tre overgangar skil seg ut med $\hat{J}_{\text{mat}} > 0{,}72$. Kvar av dei samanfell med ei kjend materiell eller teknologisk omvelting.

\textbf{1. Jugend til funksjonalisme ($\hat{J}_{\text{mat}} = 0{,}853$).} Dette er det st\o rste landskapsskiftet i heile datasettet. Materialpaletten vert n\ae rast fullstendig utskifta: jugend sin handverkstradisjon i tre og smijarn vert erstatta av funksjonalismen sitt repertoar av st\r{a}alr\o yr, kryssfiner, plast og aluminium. Teknikkrepertoaret endrar seg tilsvarande ($\hat{J}_{\text{tekn}} = 0{,}750$): formboyging, sveising og spr\o ytest\o yping erstattar tapping og utskjering. Materialentropien stig med $\Delta H' = 1{,}53$ bits, det st\o rste einskildspranget mellom to stilar. Denne overgangen korresponderer med det \citet{giedion1948} kalla ``mekaniseringa'', men dataa viser at det ikkje er ein gradvis overgang: det er eit diskret landskapsskifte der heile topografien vert rekonfigurert.

\textbf{2. Historisme til jugend ($\hat{J}_{\text{mat}} = 0{,}741$).} Historismen sin eklektiske materialpalett, driven av industriell masseproduksjon, vert erstatta av jugend sin reaksjon\ae re handverksestetikk. Materialentropien fell med $\Delta H' = 0{,}98$ bits: jugend \emph{snevrar inn} det materielle repertoaret medvite, som ein estetisk programerkl\ae ring. Denne overgangen er interessant fordi han g\r{a}ar \emph{mot} den generelle entropitrenden: det er eit landskapsskifte drive av kultur, ikkje av teknologi.

\textbf{3. Empire til historisme ($\hat{J}_{\text{mat}} = 0{,}724$).} Empire sin konsentrerte mahognipalett vert erstatta av historismen si eksplosive materialmengd: 64 stolar i 24 ulike materialar. Dette er industrialiseringa i materiell form: dampmaskiner, nye limtypar, import av eksotiske tresortar fr\r{a}a nye kjelder. Materialentropien stig med $\Delta H' = 0{,}77$ bits.


\subsection{Stabile overgangar og kontinuitet}

Ikkje alle stilskifte er landskapsskifte. Overgangen fr\r{a}a nyklassisisme til empire ($\hat{J}_{\text{mat}} = 0{,}421$) er den mest stabile i heile serien: stilane deler storparten av materialpaletten (mahogni, eik, silke, polstring). Empire er ikkje eit nytt landskap; det er ein ny posisjon i \emph{same} landskap. Stilskiftet er estetisk, ikkje materiell.

Tilsvarande er overgangen fr\r{a}a barokk til r\'{e}gence ($\hat{J}_{\text{mat}} = 0{,}481$) relativt stabil: begge stilane opererer innanfor det koloniale materialsystemet der eik, mahogni og polstringsmaterialar dominerer. R\'{e}gence er ein raffinert variant av barokkens materialpalett, ikkje eit brot med han.

Monsteret er tydeleg: \textbf{dei storste Jaccard-avstandane korresponderer med materielle revolusjonar, ikkje med estetiske revolusjonar.} Empire til historisme er industrialiseringa. Historisme til jugend er handverksreaksjonen. Jugend til funksjonalisme er st\r{a}al- og plastrevolusjonen. Stilskifte som \emph{ikkje} involverer materiell omvelting (barokk til r\'{e}gence, nyklassisisme til empire) er navigasjon innanfor same landskap: posisjonsendring, ikkje terrengendring.


% ================================================================
\section{Diskusjon}

\subsection{Fr\r{a}a metafor til m\r{a}aling}

Denne artikkelen har gjort det \citet{wright1932} og \citet{kauffman1993} ikkje gjorde: fylt fitnesslandskapet med empiriske verdiar. Kvar av dei 13 stilperiodane har no ein m\r{a}alt posisjon i eit femdimensjonalt rom. Landskapet er ikkje lenger ei skisse p\r{a}a ein tavle; det er ein tabell med tal.

Tre hovudfunn fortener framheving:

\textbf{Fyrst:} Stilperiodar er ikkje punkt i formrommet; dei er \emph{regionar} med distinkt topografi. Empire er ein n\r{a}al ($\sigma_D = 0{,}066$); postmodernisme er ein vulkan ($\sigma_D = 0{,}351$); rokokko er ein internasjonal rygg ($G_s = 1{,}82$). Landskapsmetaforen er ikkje berre poetisk; han er kvantitativt presis.

\textbf{Dernest:} Landskapsskifte er materielle, ikkje estetiske. Dei tre st\o rste Jaccard-avstandane ($\hat{J} > 0{,}72$) korresponderer alle med teknologiske omveltingar: industrialisering, handverksreaksjon, st\r{a}alrevolusjon. Estetiske stilskifte (barokk til r\'{e}gence, nyklassisisme til empire) er navigasjon innanfor same terreng. Dette stadfestar FFF sin kjernep\r{a}astand: det er \emph{materialet} som definerer landskapet, ikkje den estetiske vilja.

\textbf{Til sist:} Intensitetsindeksen avdekkjer eit overraskande moenster. Det mest ekstreme toppunktet er ikkje barokken (som ein kanskje ville forvente fr\r{a}a konvensjonell designhistorie), men postmodernismen ($\mathcal{I} = 182{,}6$). Postmodernismen investerer meir materiale per stol enn nokon annan stilperiode, samstundes som han varierer mest. Det er det dyraste oppr\o ret i designhistoria. Barokken ($\mathcal{I} = 129{,}2$) er meir fokusert: h\o g intensitet kanalisert gjennom eit smalare formrom.


\subsection{Kvifor rugleheit?}

\citet{kauffman1993} viste at i $NK$-landskap aukar talet p\r{a}a lokale optimum eksponentielt med koplingsgraden $K$. N\r{a}ar komponentane i eit system interagerer sterkt, vert global optimalisering umogleg, og systemet fangar seg i lokale toppunkt.

Me observerer 13 distinkte toppunkt (stilperiodar) i eit fem\-dimen\-sjonalt landskap. Dette er konsistent med eit moderat ruglete terreng: dei fem dimensjonane interagerer (materialval p\r{a}averkar teknikkrepertoar, teknikk p\r{a}averkar dimensjonar, geografi p\r{a}averkar materialtilgang), men ikkje s\r{a}a sterkt at landskapet vert kaotisk. \citet{gavrilets2004} sine ``holey landscapes'' tilf\o yer eit viktig perspektiv: mellom dei lokale toppunkta kan det finst n\o ytrale r\aa s (ridges) der eit system kan forflytte seg fr\r{a}a eitt toppunkt til eit anna utan \r{a}a g\r{a}a gjennom ein dal. Dei stabile stilovergangane i tabell~\ref{tab:jaccard} (nyklassisisme til empire, $\hat{J} = 0{,}421$) kan representere slike n\o ytrale r\aa s: eit stilskifte utan landskapsskifte.


\subsection{Avgrensingar}

Fire atterhald er naudsynte:

(1) Berre 13 stilperiodar passerer kravet om $N \geq 17$. Finare stilinndelingar (chippendale, sheraton, dragestil, biedermeier) har for f\r{a}a observasjonar til statistisk analyse. Ein st\o rre database ville tillate h\o gare oppl\o ysing i landskapet.

(2) Dei fem dimensjonane er proksivar: dei m\r{a}aler ikkje seleksjonstrykka direkte, men indirekte gjennom deira effektar p\r{a}a form, materiale og teknikk. Ein ideell studie ville m\r{a}ale seleksjonstrykka sjolve (produksjonskostnad, kulturell aksept, ergonomisk yting).

(3) Museumssamlingar overrepresenterer prestisjeobjekt. Den gjennomsnittlege empirestolen i museet er truleg meir ``typisk'' empire enn den gjennomsnittlege empirestolen som faktisk vart produsert. Denne biasen kan gi kunstig l\r{a}ag $\sigma_D$ for nokon stilar.

(4) Jaccard-avstanden behandlar alle materialsubstitusjonar som likeverdige: \r{a}a erstatte eik med mahogni tel like mykje som \r{a}a erstatte eik med plast. Ein vekta avstandsm\r{a}al som reflekterer materiala sine affordanseavstandar ville gje eit meir nyansert bilete av landskapsskifta.


% ================================================================
\section{Konklusjon}

Fitnesslandskapet for europeisk stoldesign er no kartlagt. 13 stilperiodar okkuperer distinkte posisjonar i eit femdimensjonalt rom definert av materialentropi, teknikkompleksitet, dimensjonell varians, geografisk spreiing og materiell intensitet.

Postmodernismen er det mest ekstreme toppunktet ($\mathcal{I} = 182{,}6$), driven av det h\o gaste materielle investeringsniv\r{a}aet i datasettet. Barokken er det nest h\o gaste ($\mathcal{I} = 129{,}2$), med meir fokusert intensitet. Empire er det strammaste ($\sigma_D = 0{,}066$): ein n\r{a}al i landskapet der materialet (mahogni), teknikken (tapping) og proporsjonslara (nyklassisisme) konvergerer mot eit punkt. Modernismen er det breiaste ($\mathcal{B} = 32{,}1$): eit internasjonalt plat\r{a}a der mange materialar, teknikkar og land er representerte.

Tre landskapsskifte, definert som $\hat{J}_{\text{mat}} > 0{,}72$, korresponderer med tre materielle revolusjonar: industrialiseringa (empire til historisme), handverksreaksjonen (historisme til jugend) og st\r{a}al-plast-revolusjonen (jugend til funksjonalisme). Stilskifte som ikkje involverer materiell omvelting har l\r{a}agare Jaccard-avstandar: dei er navigasjon innanfor same terreng, ikkje terrengendring.

Fitnesslandskapet er ikkje ein metafor. Det er eit kart med koordinatar, toppunkt og overgangar. Designhistoria er ikkje ein monoton marsj mot det optimale; det er navigasjon over ruglete terreng der kvar topp er midlertidig, kvar dal er ein overgang, og kvar materiell revolusjon er eit nytt landskap.

\begin{quote}
\textbf{Fitnesskartet erstattar metaforen.}
\end{quote}


% ================================================================
\newpage
\onecolumn
\bibliographystyle{plainnat}
\bibliography{references}

\end{document}
